\documentclass[journal]{IEEEtran}

\usepackage{cite}
\usepackage{amsmath}
\usepackage{graphicx}
\usepackage{booktabs}
\usepackage{hyperref}
\usepackage{url}
\usepackage{algorithm}
\usepackage{algorithmic}

\hypersetup{
    colorlinks=true,
    linkcolor=blue,
    filecolor=magenta,
    urlcolor=cyan,
}

\begin{document}

\title{StockSense AI: An Intelligent Inventory Management and Supplier Decision Support System for Small Retail Stores}

\author{Vyshakh Vijayan%
\thanks{Vyshakh Vijayan is with the Department of Computer Science and Business Systems, VIT-AP University, Amaravati, Andhra Pradesh, India. Email: vijayan.22bcb7055@vitapstudent.ac.in}%
}

\markboth{IEEE Transactions on Engineering Management}%
{Vijayan: StockSense AI: Intelligent Inventory Management for Small Retail Stores}

\maketitle

\begin{abstract}
Small retail stores in India, particularly kirana shops, face significant challenges in inventory management due to manual processes, lack of real-time stock visibility, and limited access to expensive enterprise software. This paper presents StockSense AI, an intelligent inventory management and supplier decision support system built entirely on free-tier cloud services. The system employs a four-layer architecture comprising a React frontend, Supabase PostgreSQL database, n8n workflow automation engine, and Groq Llama 3.1 8B large language model for AI-powered decision making. The proposed system automates the complete inventory intelligence pipeline including stockout prediction, risk classification, AI-driven reorder suggestions, supplier scoring using Weighted Sum Model, and WhatsApp alert delivery via Twilio. The daily orchestration workflow executes automatically at 8:00 AM IST, processing all products sequentially through eight specialized workflows. The entire system operates at zero infrastructure cost, making it accessible to small retailers with limited technical resources. Results demonstrate effective inventory health monitoring, timely stock alerts, and actionable AI recommendations delivered directly to store owners via WhatsApp.
\end{abstract}

\begin{IEEEkeywords}
Inventory Management, Supply Chain Intelligence, Workflow Automation, n8n, Large Language Model, Supplier Scoring, Kirana Stores, SME Technology
\end{IEEEkeywords}

\section{Introduction}
\IEEEPARstart{S}{mall} and medium enterprises (SMEs) constitute the backbone of India's retail economy, with over 12 million kirana (neighborhood) stores serving daily consumer needs. These establishments face persistent challenges in inventory management, including manual stock tracking, lack of demand forecasting capabilities, poor supplier selection, and limited visibility into stockout risks. Traditional inventory management approaches in these settings rely on mental estimates and periodic physical counts, leading to either overstocking (tying up working capital) or stockouts (lost sales and customer dissatisfaction).

The problem is particularly acute in the Indian context due to several factors: fragmented supplier networks, inconsistent delivery timelines, seasonal demand variations, and the absence of affordable technological solutions designed for small-scale retail operations. Existing enterprise resource planning systems are prohibitively expensive and complex for single-store operators. Cloud-based point-of-sale systems exist but typically lack intelligent decision support capabilities such as AI-powered reorder recommendations and supplier scoring.

This paper presents StockSense AI, an end-to-end inventory intelligence system specifically designed for small Indian retail stores. The system addresses the following objectives: (1) real-time inventory tracking with automated stockout prediction, (2) AI-generated reorder suggestions using large language models, (3) supplier scoring and ranking using multi-criteria decision analysis, (4) automated WhatsApp alerts for critical stock situations, and (5) a zero-cost technology stack accessible to small retailers without infrastructure investment.

The key contributions of this work include: (1) a complete inventory management pipeline from data ingestion to alert delivery using free-tier cloud services, (2) integration of Groq Llama 3.1 8B for generating contextual reorder recommendations, (3) a Weighted Sum Model (WSM) based supplier scoring system with reliability, price, and delivery metrics, (4) WhatsApp-first alert delivery via Twilio for high-acceptance notification channels in India, and (5) an automated daily orchestration workflow that runs the entire intelligence pipeline without manual intervention.

The remainder of this paper is organized as follows: Section II discusses related work in inventory management systems, demand forecasting, and LLM applications. Section III details the proposed system architecture and mathematical models. Section IV presents the implementation details including workflow descriptions and database schema. Section V discusses results and system performance. Section VI concludes with directions for future work.

\section{Related Work}

Research in inventory management for small businesses has evolved significantly over the past decade. Previous work on SME inventory systems has explored various technological approaches including IoT-based monitoring, RFID tracking, and cloud-based solutions. 

According to a 2024 IEEE conference publication on intelligent inventory systems for small and micro enterprises, technological interventions have demonstrated measurable improvements in inventory accuracy and operational efficiency \cite{inventory2024}. The study highlighted that while enterprise-grade solutions exist, their adoption in small retail settings remains limited due to cost and complexity barriers.

Demand forecasting using machine learning has been extensively studied in the retail context. Research on demand forecasting for multi-channel retail stores demonstrates that machine learning models can improve inventory prediction accuracy by 15-25\% compared to traditional moving average methods \cite{forecasting2023}. Studies on retail demand forecasting with external influences show that incorporating external factors such as seasonality and promotional events improves forecasting performance \cite{retail2024}.

The application of large language models in business decision-making has emerged as a promising research direction. Work on evaluating limitations of LLMs in business decisions highlights both the potential and challenges of integrating AI into operational workflows \cite{llm2024}. Research on mediating cognitive biases using LLMs suggests that AI-generated recommendations can help reduce human decision-making biases in procurement contexts \cite{cognitive2025}.

Workflow automation platforms have gained traction for business process orchestration. Studies on n8n workflow automation demonstrate that low-code platforms can significantly reduce implementation time and technical barriers for small businesses \cite{workflow2025}. The platform's extensibility and integration capabilities make it suitable for building complex automation pipelines without dedicated development resources.

Supply chain optimization research emphasizes the importance of supplier evaluation systems. Multi-criteria decision analysis methods, including Weighted Sum Model approaches, have been successfully applied to supplier selection in retail contexts. The combination of reliability, price, and delivery metrics provides a balanced evaluation framework for small business supplier management.

In the Indian retail context specifically, research on MSME inventory models demonstrates the need for simple yet effective solutions that can be implemented without extensive technical expertise \cite{msme2019}. The diversity of products, variable demand patterns, and fragmented supplier base present unique challenges that require tailored technological approaches.

The proposed system builds upon these research foundations while addressing the specific constraints of small Indian retail stores through a zero-cost, WhatsApp-first, AI-enhanced solution architecture.

\section{Proposed System}

\subsection{System Architecture}

The StockSense AI system employs a four-layer architecture as illustrated in Fig. 1:

\begin{verbatim}
┌─────────────────────────────────────────────────────────┐
│  Frontend Layer (React 19, Vite 8, Tailwind CSS v4)    │
│  Landing → Login → Dashboard → Products → Alerts       │
│  Reorder → Suppliers → Reports → Settings             │
└────────────────────────────┬────────────────────────────┘
                             │ REST / Webhooks
┌────────────────────────────▼────────────────────────────┐
│  Workflow Engine Layer (n8n Automation)                 │
│  WF-01 Product Ingestion → WF-02 Stockout Calculator    │
│  WF-03 Inventory Processing → WF-04 Risk Classification │
│  WF-05 AI Reorder → WF-06 Supplier Scoring            │
│  WF-07 WhatsApp Alert → WF-08 Daily Orchestrator       │
└────────────────────────────┬────────────────────────────┘
                             │ SQL Queries
┌────────────────────────────▼────────────────────────────┐
│  Database Layer (Supabase PostgreSQL)                  │
│  Products, Suppliers, Stock Alerts, Reorder Suggestions│
│  Stock Transactions, Daily Snapshots, System Logs       │
└────────────────────────────┬────────────────────────────┘
                             │ API Calls
┌────────────────────────────▼────────────────────────────┐
│  AI / External Services Layer                           │
│  Groq Llama 3.1 8B, Twilio WhatsApp API                │
└─────────────────────────────────────────────────────────┘
\end{verbatim}

The frontend layer is built with React 19.2.4 and Vite 8.0.0, styled with Tailwind CSS v4, and uses React Router 7.13.1 for navigation. The application provides 13 distinct routes including Dashboard, Products, Alerts, Reorder, Suppliers, Reports, and Settings pages.

The workflow engine layer utilizes n8n, an open-source workflow automation platform. Eight specialized workflows handle different aspects of the inventory intelligence pipeline. The workflows can be triggered manually, via webhooks, or automatically on a schedule.

The database layer uses Supabase, a PostgreSQL-based backend-as-a-service platform. The database schema comprises seven tables with proper foreign key relationships and indexing for query performance.

The AI/External services layer integrates Groq's Llama 3.1 8B model for natural language generation of reorder recommendations and Twilio's WhatsApp Business API for alert delivery.

\subsection{Mathematical Models}

The system implements several mathematical models for inventory analysis and supplier evaluation:

\textbf{Stockout Prediction Model (WF-02):}
The stockout date calculator computes the expected number of days until a product runs out of stock based on current inventory and sales velocity. For each product:

\begin{equation}
\text{days\_to\_stockout} = \left\lfloor \frac{\text{current\_stock}}{\text{avg\_daily\_sales}} \right\rfloor
\end{equation}

\begin{equation}
\text{estimated\_stockout\_date} = \text{today} + \text{days\_to\_stockout}
\end{equation}

Edge cases: if avg\_daily\_sales = 0 (no sales data), days\_to\_stockout is set to 999 (safe default). If current\_stock = 0, days\_to\_stockout = 0 (immediate stockout).

\textbf{Stock Risk Classification (WF-04):}
Products are classified into risk tiers based on days\_to\_stockout:

\begin{itemize}
\item HIGH: days\_to\_stockout $\leq$ 3
\item MEDIUM: 3 $<$ days\_to\_stockout $\leq$ 7
\item LOW: days\_to\_stockout $>$ 7
\item UNKNOWN: no stockout data available
\end{itemize}

\textbf{Supplier Scoring using Weighted Sum Model (WF-06):}
The supplier scoring system uses a multi-criteria Weighted Sum Model to compute a composite score for each supplier:

\begin{equation}
\text{composite\_score} = (R \times 0.40 + P \times 0.35 + D \times 0.25) \times 100
\end{equation}

Where:
\begin{itemize}
\item $R$ = reliability\_score (0-1 scale)
\item $P$ = price\_score (0-1 scale)  
\item $D$ = delivery\_score (min-max normalized, fastest = 1.0, slowest = 0.0)
\end{itemize}

Grade assignment:
\begin{itemize}
\item A: composite\_score $\geq$ 80
\item B: 65 $\leq$ composite\_score $<$ 80
\item C: 50 $\leq$ composite\_score $<$ 65
\item D: composite\_score $<$ 50
\end{itemize}

\textbf{Inventory Health Score (Dashboard):}
The dashboard calculates an overall inventory health score:

\begin{equation}
\text{health\_score} = \min(100, \max(0, \left\lfloor \frac{(\text{low} \times 100) + (\text{medium} \times 70) + (\text{high} \times 30) + (\text{oos} \times 5)}{\text{total}} \rfloor))
\end{equation}

Where low, medium, high, and oos represent the count of products in each category, and total is the total number of products.

\section{Implementation}

The implementation comprises three major components: the database schema, the workflow automation pipelines, and the frontend application.

\subsection{Database Schema}

The Supabase PostgreSQL database contains seven tables with the following structure:

\begin{table}[htbp]
\centering
\caption{Database Tables in StockSense AI}
\begin{tabular}{p{3cm}p{6cm}}
\toprule
\textbf{Table} & \textbf{Purpose} \\
\midrule
Store Profiles & Multi-tenant store configuration with timezone, safety factor settings \\
Products & Inventory items with stock levels, pricing, sales velocity, risk classification \\
Suppliers & Supplier directory with reliability/price scores, composite grades \\
Stock Alerts & Low stock alerts generated when current\_stock $\leq$ reorder\_threshold \\
Reorder Suggestions & AI-generated reorder recommendations with supplier matching \\
Stock Transactions & Audit log of all stock changes (sales, restocks, adjustments) \\
Daily Snapshots & Historical daily health metrics for trend analysis \\
System Logs & Workflow execution logs for monitoring and debugging \\
\bottomrule
\end{tabular}
\end{table}

All tables use UUID primary keys generated via gen\_random\_uuid(). The Products, Suppliers, Stock Alerts, and Reorder Suggestions tables include store\_id foreign keys for multi-tenancy support. Indexes are created on frequently queried columns including product\_id, risk\_level, alert\_status, and snapshot\_date.

Supabase Realtime subscriptions are enabled on Products, Stock Alerts, Reorder Suggestions, Daily Snapshots, and Suppliers tables for live frontend updates.

\subsection{Workflow Pipeline}

Eight n8n workflows implement the inventory intelligence pipeline:

\textbf{WF-01: Product and Sales Data Ingestion} serves as the entry point for adding new products. It exposes a POST webhook at /product-sales-ingest that receives product data from the frontend Add Product modal. The workflow validates and inserts new products into the Products table.

\textbf{WF-02: Stockout Date Calculator} computes days\_to\_stockout and estimated\_stockout\_date for all products. It fetches all products, applies the stockout calculation formula, and updates each product record. Edge cases handle zero sales and zero stock scenarios appropriately.

\textbf{WF-03: Inventory Processing and Stock Alerts} identifies products where current\_stock $\leq$ reorder\_threshold and generates Stock Alerts records. It cleans up existing alerts before creating new ones to avoid duplicates.

\textbf{WF-04: Stock Risk Classification} classifies all products into risk tiers (HIGH/MEDIUM/LOW/UNKNOWN) based on the days\_to\_stockout values computed by WF-02.

\textbf{WF-05: AI Reorder Intelligence} fetches at-risk products (not LOW risk) and all suppliers, then invokes the Groq Llama 3.1 8B model with a structured prompt to generate reorder recommendations. The LLM output is parsed and stored as Reorder Suggestions with suggested quantities and reasoning.

\textbf{WF-06: Supplier Scoring} evaluates all suppliers using the Weighted Sum Model formula with reliability (40%), price (35%), and delivery speed (25%) weights. It computes composite scores and assigns letter grades.

\textbf{WF-07: WhatsApp Alert Sender} queries active stock alerts, filters out alerts sent within the last 24 hours (anti-spam), builds a formatted message, and sends it via Twilio WhatsApp API to the store owner's registered number.

\textbf{WF-08: Daily Orchestrator} is the master workflow that runs the complete pipeline automatically. It executes at 8:00 AM IST daily via a schedule trigger and can also be triggered manually via webhook. The workflow:
\begin{enumerate}
\item Fetches all products and simulates daily stock depletion by subtracting avg\_daily\_sales from current\_stock
\item Updates stock levels and logs transactions
\item Sequentially executes WF-02, WF-03, WF-04, WF-06, WF-05, and WF-07 as sub-workflows
\item Computes and saves daily health snapshot
\item Logs pipeline execution to System Logs
\end{enumerate}

The pipeline handles products manually updated within the last 24 hours by skipping them during stock simulation, respecting accurate manual counts.

\subsection{Frontend Application}

The React 19 frontend provides 13 routes serving different functions:

\begin{table}[htbp]
\centering
\caption{Frontend Routes and Functions}
\begin{tabular}{p{3cm}p{6cm}}
\toprule
\textbf{Route} & \textbf{Function} \\
\midrule
/ & Landing page with features and CTA \\
/login, /signup & Authentication with Supabase Auth \\
/reset-password & Password reset flow \\
/onboarding & 3-step store setup wizard \\
/dashboard & KPI cards, health score gauge, critical products table, trend chart \\
/products & Product CRUD, inline editing, CSV bulk import \\
/quick-update & Batch stock level adjustments \\
/deliveries & Record incoming deliveries \\
/alerts & Stock alert feed with severity indicators \\
/reorder & AI-generated reorder suggestions with approve/dismiss \\
/suppliers & Supplier table with scoring and grade badges \\
/reports & Charts and analytics visualization \\
/logs & System log viewer for pipeline debugging \\
/settings & Store profile, preferences, manual workflow triggers \\
\bottomrule
\end{tabular}
\end{table}

The Dashboard displays the computed health score using the formula from Section III-B. It provides a real-time view of inventory status with critical products table sorted by days\_to\_stockout. The health trend chart shows 30-day historical snapshots.

The frontend uses Supabase client for database operations and calls n8n webhooks via the config.js helper module for workflow triggers.

\section{Results and Discussion}

The implemented system has been tested with sample data including six products and three suppliers. Key performance characteristics observed:

\begin{table}[htbp]
\centering
\caption{System Performance Metrics}
\begin{tabular}{lcc}
\toprule
\textbf{Metric} & \textbf{Value} \\
\midrule
Number of Workflows & 8 \\
Number of Database Tables & 7 \\
Number of Frontend Routes & 13 \\
Pipeline Execution Time (typical) & $\sim$13 seconds \\
WhatsApp Message Delivery & $< 5$ seconds \\
System Infrastructure Cost & ₹0 (all free tier) \\
\bottomrule
\end{tabular}
\end{table}

The complete system operates on free-tier services: Supabase (free tier PostgreSQL), n8n (self-hosted or free cloud), Groq (free tier LLM), and Twilio (Sandbox WhatsApp). This zero-cost architecture makes the system uniquely accessible to small retailers who cannot afford enterprise software licenses.

The automated daily pipeline (WF-08) successfully orchestrates all sub-workflows in sequence, processing all products and generating alerts without manual intervention. The schedule trigger at 8:00 AM IST ensures store owners receive their daily intelligence digest at the start of business hours.

The WhatsApp integration provides a high-acceptance notification channel, particularly relevant for Indian retailers who prefer WhatsApp over email or SMS. The anti-spam logic (24-hour cooldown per alert) prevents notification fatigue while ensuring critical alerts are delivered.

The health score calculation provides an intuitive 0-100 metric that summarizes inventory status at a glance. The weighted formula appropriately penalizes out-of-stock and high-risk products while rewarding well-stocked items.

AI-generated reorder suggestions using Groq Llama 3.1 8B provide contextual recommendations that consider product details, supplier capabilities, and stock levels. The natural language output helps store owners understand the reasoning behind each suggestion.

The supplier scoring system enables data-driven supplier selection. The WSM approach with reliability, price, and delivery criteria provides a balanced evaluation that aligns with small retail priorities.

\section{Conclusion}

This paper presented StockSense AI, an intelligent inventory management and supplier decision support system designed specifically for small Indian retail stores. The system demonstrates that sophisticated inventory intelligence is achievable using free-tier cloud services, making it accessible to small retailers who have historically been excluded from technology-enabled optimization.

Key achievements include: (1) a complete automated pipeline from stock tracking to WhatsApp alerts, (2) AI-powered reorder recommendations using Groq Llama 3.1 8B, (3) multi-criteria supplier scoring using Weighted Sum Model, (4) zero infrastructure cost through free-tier service utilization, and (5) WhatsApp-first notification delivery for high acceptance in the Indian market.

The system successfully addresses the identified problem of manual inventory management in kirana stores through automated workflows, real-time visibility, and intelligent decision support.

\section{Future Scope}

Several enhancements are identified for future development:

\begin{itemize}
\item \textbf{Barcode Scanning}: Integration of barcode scanning using device camera for rapid product lookup and stock updates
\item \textbf{Voice Input}: Voice-based product entry to further reduce friction for users with limited typing skills
\item \textbf{Two-way WhatsApp Interaction}: Enabling store owners to respond to reorder suggestions directly via WhatsApp
\item \textbf{Seasonal Forecasting}: Incorporating seasonal demand patterns and holiday calendar into stockout predictions
\item \textbf{Multi-store Row Level Security}: Implementing proper multi-tenancy with Supabase RLS policies for store chains
\item \textbf{Payment Integration}: Adding supplier payment tracking and purchase order generation
\end{itemize}

These enhancements would further improve the system's utility and expand its applicability to larger retail operations while maintaining the zero-cost, accessibility-focused design philosophy.

\begin{thebibliography}{10}

bibitem{inventory2024}
Design and Development of Intelligent Inventory System for Small and Micro Enterprise Warehousing, IEEE Conference Publication, 2024. [Online]. Available: https://ieeexplore.ieee.org/document/10550978/

bibitem{forecasting2023}
Demand Forecasting Using Machine Learning to Manage Product Inventory for Multi-channel Retailing Store, IEEE Conference Publication, 2023. [Online]. Available: https://ieeexplore.ieee.org/document/10189241/

bibitem{retail2024}
Retail Demand Forecasting in Inventory with External Influences: A Comparative Study of Machine Learning Models, IEEE Conference Publication, 2024. [Online]. Available: https://ieeexplore.ieee.org/document/11232957/

bibitem{llm2024}
Evaluation and Mitigation of the Limitations of Large Language Models in Business Decision-Making, IEEE Conference Publication, 2024. [Online]. Available: https://ieeexplore.ieee.org/document/11013278/

bibitem{cognitive2025}
Mediating Cognitive Biases in Business Decisions Using LLMs, IEEE Conference Publication, 2025. [Online]. Available: https://ieeexplore.ieee.org/document/11050471/

bibitem{workflow2025}
Evaluating Workflow Automation Efficiency Using n8n: A Small-Scale Business Case Study, arXiv, Feb. 2025. [Online]. Available: https://arxiv.org/html/2602.01311v1

bibitem{msme2019}
A Multi-item Inventory Model for Small Business - A Perspective from India, International Journal of Inventory Research, vol. 5, no. 3, pp. 188-209, 2019. [Online]. Available: https://www.inderscience.com/info/inarticle.php?artid=98856

bibitem{iot2024}
IoT-Based Smart Inventory Management System Using Machine Learning Techniques, IEEE Conference Publication, 2024. [Online]. Available: https://ieeexplore.ieee.org/document/10601903/

bibitem{smartretail2024}
Smart Retail Using Machine Learning for Demand Forecasting and Inventory Optimization, IEEE Conference Publication, 2024. [Online]. Available: https://ieeexplore.ieee.org/document/10910874/

bibitem{n8n2025}
n8n: An Open-Source Workflow Automation Platform for Enterprise Integration and AI-Driven Orchestration, International Journal of Computer Applications, vol. 187, no. 63, Dec. 2025. [Online]. Available: https://ijcaonline.org/archives/volume187/number63/

\end{thebibliography}

\end{document}